\section{案例研究：机器人控制}

\subsection{Offline policy in robotics}

在 robotics 任务中，数据采集成本高，offline RL 可大幅减少 real-world trials。典型做法：

\begin{itemize}
    \item 在 simulator 训练 policy 并在 physical robot 上 replay 关键 traj。
    \item 用 offline dataset 中的 reward multiplier 强调安全相关动作（如 low velocity）。
    \item 记录 each deployment run 中 sensor readings 以便对错误进行 counterfactual audit。
\end{itemize}

\begin{table}[H]
\centering
\begin{tabular}{p{0.25\textwidth} p{0.2\textwidth} p{0.2\textwidth} p{0.2\textwidth}}
\toprule
Task & Data source & Safety guardrail & Offline method \\
\midrule
Arm manipulation & Robot logs + human demos & Collision penalties & CQL \\
Drone navigation & Simulation + recorded flights & Wind model alert & IQL + replay \\
\bottomrule
\end{tabular}
\caption{机器人场景中的 offline RL 配置}
\end{table}

\subsection{本章小结}

机器人案例强调：data source 多元、reward carefully shaped、monitoring & replay 必不可少。

